\section{4. Method}

\subsection{4.1 Overview}

BAM-UMoE augments attention with an MoE refinement of the value pathway. 
The forward pass:

1. \textbf{Standard QKV projections} with frozen base weights ( 4.2).
2. \textbf{MoE value refinement} applied additively to V ( 4.3).
3. \textbf{Gated scaling} for controlled ramp-up ( 4.4).
4. \textbf{Output projection} with nested bypass ( 4.5).

\textbf{Figure 4.1} shows the architecture.

\begin{lstlisting}
    +------------------------------------------+
    |  Input X                                  |
    +------------------+-----------------------+
                       |
        +--------------+--------------+
        |              |              |
                                       
    +-------+      +-------+      +-------+
    | Q     |      | K     |      | V     |   $\leftarrow$ frozen base
    | proj  |      | proj  |      | proj  |   + LoRA
    +---+---+      +---+---+      +---+---+
        |              |              |
        |              |               
        |              |      +--------------+
        |              |      |  MoE Branch  |
        |              |      |  (8 experts, |
        |              |      |   top-2)     |
        |              |      +------+-------+
        |              |             |
        |              |                $\times$ 0.1 $\times$ $\sigma$(gate)
        |              |        +---------+
        |              |        | Add     |
        |              |        | to V    |
        |              |        +----+----+
        |              |             |
        +--------------+-------------+
                       |
                        
              +----------------+
              |  Attention     |
              |  (softmax)     |
              +--------+-------+
                       |
                        
              +----------------+
              | Output proj    |
              +--------+-------+
                       |
                        
                      Y
\end{lstlisting}

\subsection{4.2 Phase 1: Base Attention}

For each projection $P \in \{Q, K, V, O\}$, we use a \texttt{LoRALinear} module:

\[
\text{Proj}_P(X) = X W_P^\top + (X A_P^\top) B_P^\top \cdot s + b_P
\tag{4.1}
\]

where $W_P$ is the frozen base, $(A_P, B_P)$ are LoRA factors, and 
$s = \alpha/r_L$ is the LoRA scaling.

Base weights are copied exactly from the pretrained attention via 
\texttt{copy\\_weights} (see Section 5).

\subsection{4.3 Phase 2: MoE Value Refinement}

\subsubsection{4.3.1 MoE Architecture}

The MoE branch produces a refinement signal of shape 
$(B, L, H \cdot d_h)$:

\[
\text{MoE}(X) = \text{TopK-Gate}(X) \cdot \text{Experts}(X)
\tag{4.2}
\]

Concretely:

\[
\begin{aligned}
X_\text{norm} &= \text{LayerNorm}(X) \\
\logits &= \text{Router}(X_\text{norm}) \in \mathbb{R}^{B \times L \times E} \\
(\text{topk\_vals}, \text{topk\_idx}) &= \text{topk}(\logits, k) \\
\text{scores} &= \text{softmax}(\text{topk\_vals}) \\
E_\text{out} &= \text{stack}([\text{Expert}_e(X) \text{ for } e \in [E]]) \\
\text{MoE}(X) &= \sum_{j=1}^{k} \text{scores}_j \cdot E_\text{out}[\text{topk\_idx}_j]
\end{aligned}
\tag{4.3}
\]

\subsubsection{4.3.2 Expert Architecture}

Each expert is a 2-layer MLP:

\[
\text{Expert}_e(X) = W_2^{(e)} \cdot \text{GELU}(W_1^{(e)} X + b_1^{(e)}) + b_2^{(e)}
\tag{4.4}
\]

with $W_1^{(e)} \in \mathbb{R}^{(2d) \times d}$ and 
$W_2^{(e)} \in \mathbb{R}^{d \times (2d)}$.

\subsubsection{4.3.3 Step-0 Zero Initialization}

\textbf{Critical.} We initialize:

\[
W_2^{(e)} = 0, \quad b_2^{(e)} = 0 \quad \forall e
\tag{4.5}
\]

This ensures $\text{Expert}_e(X) = 0$ for all inputs, so 
$\text{MoE}(X) = 0$ at step 0.

\subsubsection{4.3.4 Auxiliary Load Balancing Loss}

We add the Switch/GShard load-balancing loss:

\[
\mathcal{L}_\text{aux} = E \cdot \sum_{e=1}^{E} f_e \cdot P_e
\tag{4.6}
\]

where $f_e$ is the fraction of tokens routed to expert $e$ and $P_e$ is 
the mean router probability. This is added to the task loss with weight 
$\lambda_\text{aux} = {{AUX_WEIGHT}}$ (default 0.01).

\subsubsection{4.3.5 Value Integration}

The MoE output is reshaped and added to $V$:

\[
V' = V + \alpha \cdot \text{MoE}(X)
\tag{4.7}
\]

with $\alpha = 0.1$ (fixed). The 0.1 scaling keeps the MoE contribution 
conservative -- even after training, the MoE never dominates the base.

\subsection{4.4 Phase 3: Gated Ramp-Up}

\subsubsection{4.4.1 The Gate}

A learnable scalar gate:

\[
g_\text{gate} = \sigma(b_g), \quad b_g \in \mathbb{R}
\tag{4.8}
\]

initialized to $b_g = -3.0$, so $g_\text{gate} \approx 0.047$ at start.

The final MoE contribution:

\[
\text{MoE}_\text{final}(X) = g_\text{gate} \cdot \text{MoE}(X)
\tag{4.9}
\]

\textbf{Rationale.} Even with zero-init, the MoE branch's gradient at step 0 is 
nonzero (since backprop through $W_2^{(e)} = 0$ gives nonzero gradient for 
$W_1^{(e)}$). The gate provides \textbf{additional} attenuation to prevent early 
drift.

\subsubsection{4.4.2 Combined Output}

The final forward pass:

\[
V' = V + \alpha \cdot g_\text{gate} \cdot \text{MoE}(X)
\tag{4.10}
\]

At step 0, $W_2^{(e)} = 0 \Rightarrow \text{MoE}(X) = 0 \Rightarrow V' = V$. 
Exact preservation.

\subsection{4.5 Phase 4: Output Projection}

The output projection follows the standard pattern:

\[
Y = \text{Attn}(Q, K, V') W_O^\top + b_O
\tag{4.11}
\]

with nested encoder bypass: for RoBERTa/BERT, \texttt{out\\_lin} is not applied 
(the outer \texttt{output.dense} performs it).

\subsection{4.6 The Full Algorithm}

\textbf{Algorithm 1:} BAM-UMoE forward pass
\begin{lstlisting}
Input:  X $\in$ R^{B$\times$L$\times$d}, mask, position_embeddings
Output: Y $\in$ R^{B$\times$L$\times$d}, aux_loss

1:  Q $\leftarrow$ LoRALinear_Q(X)                          Frozen base + LoRA
2:  K $\leftarrow$ LoRALinear_K(X)
3:  V $\leftarrow$ LoRALinear_V(X)
4:
5:  if RoPE: (Q, K) $\leftarrow$ apply_rotary(Q, K, cos, sin)
6:  if GQA: K, V $\leftarrow$ repeat_kv(K, V, H/H_kv)
7:
8:  if not cross_attention:
9:      MoE_out $\leftarrow$ UnifiedMoE(X)                  Top-2 MoE
10:     V $\leftarrow$ V + 0.1   MoE_out                    MoE refinement
11:
12: S $\leftarrow$ (Q @ K^T) / $\sqrt{}$d_h                         FP32
13: S $\leftarrow$ apply_mask(S, causal, padding)
14: A $\leftarrow$ softmax(S)
15: C $\leftarrow$ A @ V
16:
17: if nested_encoder: Y $\leftarrow$ C
18: else: Y $\leftarrow$ LoRALinear_O(C)
19:
20: return Y, MoE.aux_loss
\end{lstlisting}

\subsection{4.7 Complexity Analysis}

\subsubsection{4.7.1 Time Complexity}

\begin{table}[H]
\centering
\begin{tabular}{ll}
\toprule
\textbf{Component} & \textbf{Time} \\
\midrule
QKV projections & $O(L d^2)$ \\
MoE router & $O(L d E)$ \\
Expert forward (top-2) & $O(2 L d \cdot 2 d)$ \\
Attention scores & $O(L^2 d_h)$ \\
Context & $O(L^2 d_h)$ \\
Output proj & $O(L d^2)$ \\
**Total** & $O(L^2 d_h + L d^2 + L d^2)$ \\
\bottomrule
\end{tabular}
\end{table}

The MoE branch adds $O(L d^2)$ per layer (since 2 experts $\times$ 2-layer MLP 
of width $2d$), comparable to the base attention cost.

\subsubsection{4.7.2 Parameter Count}

For \texttt{UnifiedMoE(dim=d, num\\_experts=E, output\\_dim=d)}:

\begin{table}[H]
\centering
\begin{tabular}{ll}
\toprule
\textbf{Component} & \textbf{Parameters} \\
\midrule
Router LN & $2d$ \\
Router & $d E$ \\
Each expert & $2d \cdot 2d + 2d \cdot 2d = 8d^2$ \\
Total experts & $8 d^2 E$ \\
**MoE total** & $8 d^2 E + O(d E)$ \\
\bottomrule
\end{tabular}
\end{table}

For $d = 768$, $E = 8$: \textbf{8 $\times$ 768  $\times$ 8 $\approx$ 37.7M params per layer}.

\textbf{This is enormous.} For DistilBERT (6 layers): \textbf{+226M params} on a 
66.9M model. Parameter count \textbf{increases by ~4.4$\times$}.

\textbf{BAM-UMoE is not a compression method.}